\appendixsection[TW]{Auswertung Tragwerksplanung (Eike Schling)}
\label{anlage:interview-schling}

Dieser Anhang fasst zentrale Erkenntnisse aus dem Gespräch mit Eike Schling zusammen. Im Mittelpunkt stehen tragwerksplanerische Fragen der Bauteilwiederverwendung: Welche Informationen müssen aus dem Spendergebäude erhalten bleiben? Wie lässt sich das Potenzial eines Bauteils bewerten? Welche Nachweise sind erforderlich, und wie können aus heterogenen Bauteilen neue, prüfbare Systeme entstehen?

\subsection{Tragwerkswissen muss mit dem Bauteil erhalten bleiben}

Bei Stahlbeton reichen äußere Abmessungen nicht aus, um die weitere Verwendbarkeit eines Bauteils einzuschätzen. Ebenso relevant sind die ursprüngliche Tragfunktion, Lagerung, Bewehrung, Verbindung und die Art der Zerlegung.

\begin{Blockquote}
Es muss eigentlich pro Gebäude eine gewisse Logik wie ein Handbuch geben, wie man eine Decke, eine Wand, einen Rahmen schlau auseinandernimmt, um ihn dann weiter zu benutzen. […] Die Gebäude müssen alle nach einem klaren Handbuch, das der Tragwerksplaner sich erdacht hat, zerlegt werden. Dazu muss es ein Logikbuch geben, warum sie so zerlegt wurden – das hat etwas mit der Bewehrung und dem Tragwerk zu tun. Und dann archiviert und digitalisiert werden, sodass wenn jemand ein neues Gebäude baut, er nachschauen kann: Was haben wir denn hier im Baukasten? […] Das große Problem von der Tragwerksseite ist, dass bei diesem Auseinandernehmen immer nur die geometrischen Dimensionen hinterlegt sind, aber bei Stahlbeton hat das am meisten mit der Bewehrung darin zu tun […] Das heißt, es müsste eigentlich der Originalzustand hinterlegt sein.
\end{Blockquote}

Für ein digitales Reuse-Werkzeug sollten Informationen zum Spendergebäude und zum entnommenen Bauteil miteinander verknüpft bleiben. Auf Gebäudeebene sind vor allem die Tragwerks- und Zerlegungslogik relevant. Am Bauteil sollten die ursprüngliche Einbausituation, Bewehrung, Verbindung sowie spätere Schnitte und Bearbeitungen nachvollziehbar sein.

Ein Katalog, der nur Geometrie und Materialtyp enthält, reicht für eine tragwerksplanerische Einschätzung nicht aus.

\subsection{Bauteilpotenzial statt nur Materialmenge}

Schling unterscheidet zwischen der Menge des wiederverwendeten Materials und der Qualität seiner neuen Nutzung. Eine hohe Wiederverwendungsquote kann zugleich wertvolle geometrische oder konstruktive Eigenschaften eines Bauteils mindern.

\begin{Blockquote}
Es wird immer eine Abwertung geben von der eigentlichen Möglichkeit. Aber man kann es so bewerten, dass man unterschiedliche Strategien hat und schaut, dass man nicht nur sagt, wie viel Material hat man angewendet, sondern dass man sich ein Bewertungskriterium macht, wie gut das Material im neuen Gebäude ausgenutzt wurde. Wenn man das ganze Material verwendet hat, es aber in Zwei-Meter-Stücke zerschnippelt und ganz kleine Räume damit geschaffen hat, ist das zwar gut, aber man hat das Potenzial nicht voll ausgenutzt. Das müsste in die neue Entwurfsmatrix mit einfließen. Dann hätte man das für ein anderes Gebäude, das unbedingt einen Acht-Meter-Raum braucht, viel schlauer verwenden können. Das gilt ja für alles.
\end{Blockquote}

Die Entwurfsbewertung sollte daher neben der Materialmenge auch berücksichtigen, ob nutzbare Länge, Spannweite und konstruktive Möglichkeiten erhalten bleiben. Diese Aspekte sollten zunächst getrennt dargestellt werden. Ein einzelner Gesamtwert würde die unterschiedlichen Qualitäten eines Bauteils und seiner neuen Verwendung zu stark vereinfachen.

\subsection{Nachweisbarkeit ist materialabhängig}

Die Beurteilung vorhandener Bauteile unterscheidet sich je nach Material. Stahl, Holz und Stahlbeton weisen unterschiedliche Möglichkeiten und Schwierigkeiten bei der Ermittlung ihrer Eigenschaften und beim Nachweis der verbleibenden Tragfähigkeit auf.

\begin{Blockquote}
Bei Stahl ist das Gute, das kann ohnehin recycelt werden. Da ist das Problem nicht so brennend, und man kann die Stahlgüte relativ gut wiederherstellen oder messen. […] Es hängt ganz stark an den Möglichkeiten, diese recycelten Bauteile für den neuen Bau wieder einzuschätzen. Das muss der Ingenieur nachweisen können. Bei Holz wird das schwerer, und bei Stahlbeton ist es besonders schwer. Da muss man ganz viel forschen, wie die Tragfähigkeit überhaupt nachgewiesen werden kann.
\end{Blockquote}

Das Reuse-Werkzeug muss werkstoffliches Recycling und die Wiederverwendung eines konkreten Bauteils klar unterscheiden. Bei der Bauteilwiederverwendung kommt es darauf an, welche Materialkennwerte und Angaben zur Tragfähigkeit verfügbar und belastbar sind.

Das Werkzeug sollte kenntlich machen, welche Informationen vorliegen und wo weitere Untersuchungen erforderlich sind. Der fachliche Nachweis bleibt Aufgabe der Tragwerksplanung.

\subsection{Kleine Bauteile können zu neuen Tragwerken kombiniert werden}

Kleine oder zugeschnittene Elemente müssen nicht auf geringe Spannweiten oder untergeordnete Funktionen beschränkt bleiben. Durch räumliche Anordnung, Schichtung und Materialkombination können neue Tragwirkungen entstehen.

\begin{Blockquote}
Wenn wir keine langen Bauteile haben, können wir eine größere Spannweite mit vielen kleinen erreichen. Das Problem der Zerkleinerung muss keine automatische Einschränkung in der Spannweite sein. Es kann helfen, wenn man ihnen eine räumliche Krümmung gibt, um die Tragwirkung wiederherzustellen. […] Es gibt auch andere schlaue Ansätze: Sandwich-Decken. Die sichtbaren Flächen sind Neubauteile, aber innerhalb des Sandwichs kann man wunderbar recycelte Sachen verbauen und verstecken. […] Und dann hybride Konstruktionen machen, wo man druckbeanspruchbare Bauteile wie Beton auf die Oberseite legt und zugbeanspruchbare Teile wie Stahl oder Holz auf die Unterseite.
\end{Blockquote}

Das Gespräch nennt drei Strategien: die räumliche Anordnung kleiner Elemente, die funktionale Schichtung in Sandwichsystemen und die Kombination von Materialien nach Beanspruchungsart.

Die Entwurfsassistenz sollte daher neben direkten Bauteilzuordnungen auch zusammengesetzte Systeme berücksichtigen. Solche Varianten bleiben Entwurfsideen und müssen anschließend tragwerksplanerisch geprüft werden.

\subsection{Nachweisbare Systeme als Voraussetzung der Anwendung}

Technisch mögliche Lösungen sind nur dann praktisch nutzbar, wenn übliche Ingenieurbüros sie mit vertretbarem Aufwand prüfen können. Schling betont deshalb die Bedeutung klarer Systeme und standardisierter Randbedingungen.

\begin{Blockquote}
Man muss sich Bausysteme überlegen, die sich danach relativ leicht nachweisen lassen. Wo ein normales Ingenieurbüro die Möglichkeit hat, das nachzuweisen. Wie sie sich treffen, wird zum Beispiel über eine standardisierte Stahlplatte geregelt, wo man von der gleichen Reibung und Kraftübertragung ausgehen kann. […] Wenn der planerische Aufwand und die Kosten zu hoch sind, wird wieder nicht nachhaltig gebaut. Man muss die Hürde senken, so etwas nachweisen zu können. Wie man solche Systeme entwirft, dass keine komplexen digitalen Prozesse mehr notwendig sind […]
\end{Blockquote}

Ein digitales Reuse-Werkzeug kann Systemprinzipien, Randbedingungen und offenen Prüfbedarf sichtbar machen. Bemessung und Ausführungsplanung bleiben Aufgaben der Fachplanung.

\subsection{Konfigurierbare Anschlüsse als Entwurfsidee}

Als konkrete Entwurfsidee beschreibt Schling eine neue Schnittstelle an geraden Schnittkanten von Stahlbetonbauteilen. Im Gespräch wird dieser Ansatz als „Print on Demand“ und als „USB-Stecker-Problem“ bezeichnet.

\begin{Blockquote}
\textbf{Eike Schling:} Stand heute würde ich sagen: Es gibt gerade Schnittkanten. An allen geraden Schnittkanten würde ich Löcher bohren und Stahlplatten am Rand installieren. Das ist dann eine neue, feste Verbindung im Stahlbau, die ich sofort nachweisen kann. Damit habe ich Bauteile, die anschlussfähig sind, und der Entwerfer muss sich keine Sorgen mehr um das bröselige Beton-Detail machen.

\textbf{Interviewer:} Das ist Print on Demand, wo man sagt, wir haben Bauteile.

\textbf{Eike Schling:} Ja, das ist das typische USB-Stecker-Problem. Man klickt an, welche Verbindung man gerne hätte.
\end{Blockquote}

Das Gespräch beschreibt eine mögliche konfigurierbare Anschlusslogik. Das digitale Reuse-Werkzeug soll künftig geeignete Anschlussprinzipien zur Auswahl stellen und die dafür benötigten Angaben benennen.

Befestigung, Kraftübertragung, Bemessung und Zulässigkeit sind gesondert zu prüfen. Der digitale Vorschlag dient als Grundlage für die weitere Fachplanung.
